\section{Declaração de Uso de Inteligência Artificial Generativa}
\label{sec:declaracao-ia}

Em conformidade com as diretrizes da disciplina para a entrega do trabalho final,
declaramos o uso pontual de Inteligência Artificial Generativa (ferramenta:
\textbf{Claude AI}, Anthropic) como suporte durante o desenvolvimento do código.
Abaixo detalhamos as áreas específicas onde a ferramenta foi utilizada, os prompts
empregados e a explicação do funcionamento do código gerado, demonstrando compreensão
da lógica integrada ao projeto --- e não o uso de uma "caixa preta".

\subsection{Cálculo de métricas (F1-score, precisão e recall)}
\label{sec:declaracao-ia-metricas}

\textbf{Arquivo afetado:} \texttt{scripts/eval\_metrics.py}

\begin{quote}
\itshape
``Como posso implementar uma função em Python usando NumPy para calcular a Matriz de
Confusão (Verdadeiros Positivos, Falsos Positivos e Falsos Negativos), a Precisão, o
Recall e o F1-Score comparando dois arrays booleanos 8x8 (um com a predição de
ocupação do tabuleiro e outro com o ground truth)? Preciso que o código evite o erro
de `divisão por zero'.''
\end{quote}

\textbf{Explicação do código gerado:} o trecho gerado e adaptado por nós calcula o
F1-score da detecção de ocupação do tabuleiro:
\begin{enumerate}[nosep]
  \item \textbf{Verdadeiros Positivos (TP):} comparação elemento a elemento entre a
        matriz de predição e a matriz de \textit{ground truth} (GT), \texttt{(pred == 1)
        \& (gt == 1)} --- casas de fato ocupadas que o modelo acertou.
  \item \textbf{Falsos Positivos (FP):} \texttt{(pred == 1) \& (gt == 0)} --- casas que o
        modelo marcou como ocupadas, mas o GT indicava vazias.
  \item \textbf{Falsos Negativos (FN):} \texttt{(pred == 0) \& (gt == 1)} --- peças que o
        modelo não detectou.
  \item \textbf{Fórmulas:} Precisão $= TP / (TP + FP)$; Recall $= TP / (TP + FN)$;
        F1 $= 2 \cdot (\text{Precisão} \cdot \text{Recall}) / (\text{Precisão} +
        \text{Recall})$. O código inclui um termo mínimo ($\epsilon = 10^{-9}$) somado
        ao denominador --- necessário para evitar divisão por zero no caso extremo de um
        tabuleiro inteiro ser lido como vazio, o que zeraria o denominador da precisão.
\end{enumerate}

\subsection{Visualização de resultados (bounding boxes com Matplotlib)}
\label{sec:declaracao-ia-viz}

\textbf{Arquivo afetado:} notebooks de visualização e geração de figuras do relatório.

\begin{quote}
\itshape
``Escreva um trecho de código usando Matplotlib que pegue uma imagem de um tabuleiro
de xadrez 480x480 retificado, faça um loop para dividi-lo em um grid de 8x8 (60x60
pixels por casa) e desenhe `bounding boxes' verdes nas casas onde o modelo acertou e
vermelhas nas casas onde o modelo errou. Como desenho esses retângulos vazados em cima
da imagem original?''
\end{quote}

\textbf{Explicação do código gerado:} usamos a IA para acelerar a geração de código
\textit{boilerplate} de plotagem (não de processamento de imagem), usado exclusivamente
para gerar as figuras do relatório:
\begin{enumerate}[nosep]
  \item Inicialização de figura e eixo (\texttt{fig, ax = plt.subplots()}), com a imagem
        de fundo (tabuleiro retificado) exibida via \texttt{ax.imshow()}.
  \item Loop duplo aninhado (\texttt{for i in range(8)}, \texttt{for j in range(8)})
        percorrendo o grid 8$\times$8.
  \item Coordenada do canto superior esquerdo de cada casa obtida multiplicando os
        índices pelo tamanho em pixels (\texttt{x = j * 60}, \texttt{y = i * 60}).
  \item Cada casa é desenhada com um objeto
        \texttt{matplotlib.patches.Rectangle}, recebendo posição, largura, altura, cor
        da borda e \texttt{facecolor='none'}. É esse último parâmetro que garante que a
        caixa fique vazada, sem cobrir a peça na imagem.
  \item A cor é definida dinamicamente a partir da nossa matriz de comparação
        predição~$\times$~\textit{ground truth}: verde em caso de acerto, vermelho em
        caso de erro. Cada retângulo é adicionado à figura com \texttt{ax.add\_patch()}.
\end{enumerate}
